\chapter{Conclusion and outlook}

\section{Conclusion}

In this thesis, we went over the concepts of axion as dark matter candidates, and the possible coupling of axion field gradient to half-integer nuclear spin. 
CASPEr-gradient was proposed aiming at direct detection of axion through this coupling. 
NMR technique, as the tool for spin precession measurement, was introduced to prepare the framework for experiment design and data analysis. By ramping magnet to different magnitudes, CASPEr-gradient is able to look for axion in from \SI{1}{\kHz} to \SI{600}{\MHz}. 

On one hand, CASPEr apparatus was designed to satisfy requirements for regular NMR experiments. On the other hand, it is specially adapted for precision measurement on small NMR signal, in order to cope with weak axion coupling. The Faraday room, cryostat and superconducting coil system provide a low-noise environment for SQUID as the magnetometer. %The technical details of signal generation, reception and processing are included. 
Concerning the NMR sample, we stated the reasons why \ch{^{129}Xe} was chosen for axion research. Moreover, we explained the advantages of methanol over other chemicals as the sample for system diagnostic. 

Analytical calculation and experimental measurements were performed to examine and characterize the system. We acquired knowledge of NMR signal signature, including amplitudes, linewidth and relaxation time. Methods for determining sample position with regard to gradiometer center and magnetic center were demonstrated. The magnetic field was proven to be steady for axion research. The noise of the system during operation were characterized as $\SI{2.2e-12}{}\,\Phi_0^{2}/\SI{}{\Hz}$, dominated by white noise. 

In addition to system diagnostic, we analyzed NMR signal in a scenario where excitation field is weak. The flux in the magnetometer is estimated to be proportional to the relaxation time. The stochastic nature was taken into account to correct the estimation on signal amplitude. 
A kinetic simulation was implemented to offer more evidence for analytical predictions. 
Criteria for discovering axion under certain coupling strength $\mathrm{g}_{aNN}$ were analyzed, which provided theoretical tools for exclusion on $\mathrm{g}_{aNN}$ when axion is not discovered. A 10-hr measurement was analyzed, excluding $\mathrm{g}_{aNN}$ above $10^{-5}\,\SI{}{\GeV^{-1}}$ at \SI{1.349550}{\MHz} with \SI{80}{\Hz} bandwidth. Lastly, we inferred from quantitative computation that, for the same total duration of measurements at different frequencies, we should reduce NMR linewidth at the sacrifice of detection bandwidth in one single measurement in order to scan larger parameter space in gradient coupling. 

\section{Outlook}

We are looking forward to performing axion measurement with sufficient calibrations. Right now we still lack in a few tools for characterize sample properties including $T_1$ time measurement and robust signal recording program. 

Besides, we expect progress from hyperpolarization of \ch{^{129}Xe} and CASPEr-gradient high-field phase. 